\section{Context, approach, and vocabulary}\label{sec:construct}

\subsection{The record and the question}

A US trademark application carries a description of the goods or services
the mark will cover. The description is written to secure scope of
protection: wide enough to cover the planned use, narrow enough to be
granted. It is therefore a forced, dated disclosure of what a firm expected
to sell. The USPTO publishes every application since 1884 with its
prosecution history, and this paper re-parses all 13.99 million case files
into dated events, so that each application can be followed from filing to
registration, to the sworn proof of continued use required between the fifth
and sixth year, and to cancellation or renewal.

The question is whether the language of the description, read against the
language of the rest of the industry at the time, predicts what happened to
the mark and to the firm behind it. Two properties of the language are
measured. \emph{Atypicality} is how unusual the description is for its
industry. \emph{Lead} is whether its language is ahead of or behind where the
industry's language was moving. Sections~\ref{sec:life}--\ref{sec:secevents}
relate the two to registration, to survival of the use-proof gate, to venture
funding and public listing. Section~\ref{sec:robust} reports what changes
under other choices.

\subsection{Two readers, one living backwards}

Imagine two readers of a filing's goods/services description, each having seen
a disjoint half of its industry. One has read only what the industry filed in
the five years before it. The other, living backwards as Merlin was said to,
has read only the five years after. Each is asked how surprising the wording
is. The measurement is how much they disagree: a description that startles the
first reader and not the second arrived before its industry's language moved
that way.

Both windows are anchored on the filing's own date, the 1,826 days before it
and the 1,826 days after, so no two filings made on different days share a
reference, and filings made on the same day fall in neither of the other's
windows. Inside a window every day counts the same. The measure asks whether
wording is unusual against a period, not whether it is ahead of a trend
inside that period; it has no momentum channel.

\subsection{What is measured}

The quantity is information content in Shannon's sense. Something that
happens with probability $p$ carries $-\log p$ of information when it
happens: a common thing carries little, a rare thing a lot. Here the things
are the themes a description draws on, and the probabilities are how often
the filing's own industry drew on each of them in the reference window. A
filing's score is the average information content of its themes under its
industry's frequencies, weighted by how heavily it uses each, less the same
average taken under the filing's own frequencies. It is zero for a
description composed exactly as the industry composes its descriptions and
grows as the filing puts weight where the industry puts little.

The apparatus is that of \citet{Murdock2017}, developed on Darwin's reading
notebooks, and \citet{Barron2018}, who extended it to French Revolution
parliamentary debates. They score a document against those before it and
again against those after, calling the two quantities \emph{novelty} and
\emph{transience} and their signed difference \emph{resonance}. Their object
was the life of ideas: how often an idea appeared, and how ideas stood in
relation to one another, across one reading life or one assembly's debate.
The object here is the composition of commercial offerings, which components
a product or service is described as having, and how that composition shifts
across an industry over time. The arithmetic carries over unchanged. The
interpretation does not: their document has one author and the reference
stream is that author's own reading, so a high value says something about a
mind; here a filing is scored against other firms' filings, and its author is
usually counsel drafting for scope of protection.

\paragraph{Themes, not topics.} The clusters the model fits are called
\emph{themes} throughout. In the source literature and in the name of the
method, latent Dirichlet allocation, the same objects are called topics, and
a document's topic is what it is about. A goods/services description is not
about anything in that sense; it lists components. A cluster of terms that
co-occur across descriptions (\emph{video, music, games, electronic,
audio}) is a component that many offerings share, not a subject. Theme is
the word for that. Topic is used here only in the name of the model.

\subsection{From a description to two numbers}\label{sec:definitions}

\paragraph{Counting words.} A description is reduced to a bag of words: the
text is lowercased, split into tokens of three or more letters, and counted,
with adjacent pairs counted as tokens in their own right so that
\emph{computer software} registers alongside \emph{computer} and
\emph{software}. Order is discarded. Take the filing for \textsc{Amazon.com},
lodged in the advertising and retail class on 18 April 1997:

\begin{quote}\small
computerized on line search and ordering service featuring the wholesale and
retail distribution of books, music, motion pictures, multimedia products and
computer software in the form of printed books, audiocassettes,
videocassettes, compact disks, floppy disks, CD ROMs, and direct digital
transmission
\end{quote}

\noindent Of its 45 distinct terms, 33 are in the theme model's vocabulary,
among them \emph{computerized}, \emph{search}, \emph{ordering},
\emph{ordering service}, \emph{retail}, \emph{distribution}, \emph{printed
books} and \emph{digital transmission}.

\paragraph{Grouping words into themes.} Each filing is then re-described in a
fixed, much smaller set of coordinates. Latent Dirichlet allocation is run
once over a stratified sample of the whole corpus, all 45 classes together: it
learns 50 clusters of words that tend to occur together and, for any filing,
the proportions in which its words draw on them. Nothing supervises this; the
clusters come from co-occurrence alone. The themes are shared across classes;
what is specific to a class is the reference, the average theme proportions
of that class's own filings in the window. The Amazon filing draws on six
themes at more than a fiftieth of its weight: $0.41$ on media and
entertainment goods (\emph{video, computer, electronic, music, games,
audio}), $0.15$ on computer and information services (\emph{services,
computer, software, information, data}), $0.13$ on retail store services
(\emph{store, retail, store services, featuring}), $0.10$ on printed matter
(\emph{paper, books, printed, cards, stationery}), $0.09$ on downloadable and
online software, and $0.07$ on marketing of consumer goods. The online
appendix lists every theme with its leading words and its weight in each
class.\footnote{\url{https://aporia.institute/tm-vocabulary/online-appendix/themes_T50.html}}

Fifty is coarse, roughly one theme per industry. Section~\ref{sec:resolution}
refits the model at 500 themes and separately at 50 themes within each class
and repeats the gate analysis under each; Section~\ref{sec:robust} reports
the resolution sweep and a second fit at the same resolution with a different
seed. Fifty is kept in Sections~\ref{sec:life} and~\ref{sec:secevents}
because it is the coarsest resolution at which every result holds, and
because its lead contrast is the lowest of the resolutions tried.

\paragraph{Comparing against the two references.} Write $P$ for the filing's
distribution over themes, and $Q_{\mathrm{past}}$ and $Q_{\mathrm{future}}$
for that distribution averaged over every same-class filing in the two
windows, $[d_i - W,\, d_i)$ and $(d_i,\, d_i + W]$, with $W = 1{,}826$ days.
How far $P$ sits from a reference $Q$ is the Kullback--Leibler divergence,
\[
  \mathrm{KL}(P \,\|\, Q) \;=\; \sum_{k} P_k \log \frac{P_k}{Q_k},
\]
a sum over the 50 themes, zero when the two distributions agree and growing as
the filing puts weight where the industry puts little. It is measured in nats,
the natural-log unit. Write
\[
  K^{-} = \mathrm{KL}(P \,\|\, Q_{\mathrm{past}}),
  \qquad
  K^{+} = \mathrm{KL}(P \,\|\, Q_{\mathrm{future}}),
\]
so $K^{-}$ is what the first reader feels and $K^{+}$ the second: the
filing's past-facing and future-facing surprise.

\paragraph{Taking the difference.} Amazon's filing is scored against 41{,}166
earlier filings and 146{,}532 later ones, giving $K^{-} = 1.286$ and
$K^{+} = 1.163$. The two axes used throughout are their average and their
signed difference:
\[
  \underbrace{A \;=\; \tfrac{1}{2}\left(K^{-} + K^{+}\right)}_{\text{atypicality}},
  \qquad
  \underbrace{L \;=\; K^{-} - K^{+}}_{\text{lead}} .
\]
For Amazon, $A = 1.225$ and $L = +0.123$. Against its own class those say
something the raw pair does not: its atypicality is at the 47th percentile,
entirely ordinary, while its lead is at the 94th. The industry moved toward
this description in the five years after it was filed and had not been writing
that way in the five years before. Nothing about the wording was strange; it
was early. $L$ is signed, positive for a leading filing and negative for a
lagging one, and the symbol \dkl{} is used interchangeably with it.

What ``early'' means is visible in the themes themselves. Of the six Amazon
drew on, the share of class-035 filings drawing on each, in the five years
before its filing and the five after:

\begin{center}\small
\begin{tabular}{lrrr}
\toprule
Theme & Amazon's weight & Before (\%) & After (\%) \\
\midrule
Retail store services              & $0.13$ & $42.6$ & $50.3$ \\
Computer and information services  & $0.15$ & $58.9$ & $59.9$ \\
Media and entertainment goods      & $0.41$ & $9.9$  & $11.3$ \\
Downloadable and online software   & $0.09$ & $2.3$  & $4.1$  \\
Printed matter                     & $0.10$ & $8.3$  & $7.8$  \\
Marketing of consumer goods        & $0.07$ & $10.5$ & $9.3$  \\
\bottomrule
\end{tabular}
\end{center}

\noindent A theme counts as drawn on when it carries at least a tenth of a
filing's weight. The three themes that carry most of Amazon's description,
retail, media goods, and the still-small online-software theme, all became
more common in the class over the following five years, the last nearly
doubling; the two minor themes it also touched became slightly less so. The
filing is closer to the class's future composition than to its past one, and
that, not any rare word, is its lead.

Differencing is what makes that separation possible. $K^{-}$ and $K^{+}$
correlate at $0.988$ on the software class, so entering both is close to
entering one variable twice; rotating them into an average and a difference
isolates the large common component from the small residual that carries the
timing, and leaves the two regressors nearly uncorrelated ($+0.036$ against
each other, where $K^{-}$ alone would give $+0.114$). The cost is that $L$ is
a small residual of two large quantities, about a sixth of either level's
spread, which Section~\ref{sec:robust} quantifies.

\begin{table}[h]\centering\small
\caption{The quantities, their sources, and the terms used here. \emph{Lagging}
and \emph{leading} describe the two signs of $L$; \emph{atypicality} is
unsigned and says nothing about direction.}
\label{tab:constructs}
\begin{tabular}{llll}
\toprule
Quantity & Definition & \citeauthor{Murdock2017} & Term used here \\
\midrule
Surprise against the past   & $K^{-}$                        & novelty    & past-facing surprise \\
Surprise against the future & $K^{+}$                        & transience & future-facing surprise \\
Their average               & $A = \tfrac{1}{2}(K^{-}+K^{+})$ & ---        & \textbf{atypicality} \\
Their signed difference     & $L = K^{-}-K^{+}$              & resonance  & \textbf{lead} (leading / lagging) \\
Its magnitude               & $|L|$                          & ---        & lead magnitude \\
\bottomrule
\end{tabular}
\end{table}

\paragraph{Why the words are not scored directly.} The themes are an extra
step, and the words themselves could be compared to the reference: build the
filing's distribution over the class vocabulary from the words it uses, do
the same for every filing in the window, and take the divergence. The paper
reports that scoring as a cross-check. It cannot be the primary one, because
a divergence measured on words mostly counts the length of the filing. The
divergence is an average surprise less a spread: for any two distributions,
\[
  \mathrm{KL}(P \,\|\, Q) \;=\; \underbrace{-\textstyle\sum_k P_k \log Q_k}_{\text{how unexpected these words are}}
  \;-\; \underbrace{\left(-\textstyle\sum_k P_k \log P_k\right)}_{\text{how spread out the filing is}} ,
\]
and the spread of a filing over words is bounded by how many distinct words it
has: a description using $k$ terms cannot spread over more than $k$ of them,
so its spread cannot exceed $\log k$, while the reference is spread over about
900 terms. A five-word description is charged the whole difference in width,
roughly five nats, before any of its content is considered. Measured rather
than derived: take 3{,}000 software-class filings of at least 150 words,
score each against one fixed reference, then truncate the same filings to a
random handful of words and score them again. Across a fiftyfold cut in length
the surprise term moves by $0.016$ nats and the spread term by $2.9$. Synthetic
filings whose words are drawn at random from the reference itself, maximally
ordinary by construction, score $5.72$ at three tokens and $1.65$ at four
hundred, tracking $\log(\text{distinct terms})$ to within a fiftieth of a
nat, and the average real filing scores marginally below that null at its
own length. Term-scored atypicality correlates with the log of a filing's
distinct term count at $-0.68$ in the software class.

Themes break the coupling because the 50 coordinates are always occupied.
Under term scoring the effective number of coordinates a filing spreads over
rises seventeenfold from the shortest descriptions to the longest; under
themes it varies by less than a third, and not monotonically: a three-word
fragment spreads across more themes than a long description, since the model
cannot tell which one it belongs to. The width term therefore stops carrying
length, and what is left is the part that was always about content.
Theme-scored atypicality correlates with the log term count at $-0.22$ in
the software class and $-0.10$ in advertising and retail, against $-0.68$
and $-0.62$ under terms.

\paragraph{What the themes see, and what they do not.} The score compares a
filing's theme proportions to the average proportions of its industry, so it
reads which themes a filing draws on and how heavily. It does not, in any
material way, read how a filing combines them. Taking each filing's two
leading themes and asking how often that pair occurs together relative to the
two themes' separate frequencies: how rare the themes are individually explains
70\% of the variance in atypicality, and the unusualness of the pairing adds
under two points on top of it. A filing assembled from ordinary parts in a
genuinely new configuration is therefore close to invisible to this measure.

Combinations can be scored in their own right, at two levels, and the two
levels behave differently. Both are computed on the software class, and both
are reported net of the filing's own lead and atypicality, so each is the
increment over what the construct already measures.

At the theme level, take every pair of themes a filing draws on (at a tenth
of its weight or more) and measure how much more or less often each pair
occurs together than the two themes' separate frequencies would predict, in
the class's filings of the five years before and the five after. The
negative log of that ratio is the pairing's information content; averaging
over a filing's pairs, weighted by their masses, and differencing the two
windows gives a \emph{pair lead}, positive when the class moved toward the
filing's pairings. Filings in the top fifth of pair lead fail the first
use-proof gate $2.3$ points more often than those in the bottom fifth
($t = 8.4$, unchanged by the controls), and pair lead correlates with lead
itself at only $-0.18$. Being early with an arrangement of themes is
penalised in its own right, at a rate comparable to being early with the
themes. The rarity of a filing's pairings, by contrast, carries nothing
once its atypicality is held ($-0.3$ points, $t = -1.1$).

At the word level the sign reverses. For each filing, count the share of
its adjacent word pairs that are new to the class --- the pair appearing
fewer than five times in the class's prior five years --- while both words
are established there. This is recombination: familiar parts, new
adjacency. Filings in the top fifth of that share fail the gate $6.1$
points \emph{less} often than those in the bottom fifth, and $2.8$ points
less with lead and atypicality held ($t = -10.4$). Word recombination and
theme-pair lead are uncorrelated ($r = -0.01$): they are different facts
about a filing. What the record shows, then, is not that combination is
invisible and neutral. Arranging themes the industry had not yet arranged
is punished like arriving early; coining phrases from the industry's own
words is rewarded. The construct the paper uses sees neither, and the two
run in opposite directions, so its blindness to combination is a bounded
omission rather than a hidden bias in one direction. The sections that
follow use lead and atypicality of the parts; these two measures bound what
they leave out.

\subsection{Within class, not across it}

Goods/services descriptions are drafted to secure scope of protection, not to
describe a product distinctively, and counsel reuse language from the USPTO's
own Acceptable Identification of Goods and Services Manual, from house
precedent, and from competitors' successful filings. Identical wording across
unrelated firms is therefore common, and it carries through to the scores:
17\% of scored registrations share a value with another filing in the same
class and year. Within a class, the measure picks up departure from a shared
drafting convention. Across classes it does not: a class where the Manual
supplies dense standard language shows a compressed distribution of
atypicality, and one where firms describe novel services in prose a wide one,
for reasons unrelated to how innovative either industry is. Every estimate is
therefore taken within class and year, and magnitudes compare across classes
only as signs and orderings.

The vocabulary is positional rather than causal. A leading filing was ahead of
its industry's language, which implies a trajectory without asserting the
filing caused it. Atypicality is textual, and is not trademark law's
distinctiveness doctrine, which governs registrability along the
generic-to-fanciful spectrum; the two attach to different objects, one to the
mark and one to the description of the goods it covers.

\subsection{What the construct is not called}\label{sec:neighbours}


\emph{Novelty} and \emph{transience}, in the sense of
\citet{Murdock2017} and \citet{Barron2018}, are the two KL levels
themselves: novelty is divergence from the past, transience divergence
from the future. Atypicality as used here is their average, and so is
neither of them alone; \emph{novelty} in particular is declined because it
names only the backward-looking half of a quantity that faces both ways.
\emph{Resonance} is their signed difference, and is arithmetically what
this paper calls lead. The word is avoided because it names a
mechanism --- the field resonating with an author's contribution --- that
this design cannot observe. A filing can sit where its class was heading
because it was imitated, because it and the class responded to the same
outside event, or by coincidence; the measure does not separate these,
and \emph{resonance} would assert the first.

\emph{Innovation} is a property of the product, not of the sentence
describing it: two firms selling the same thing can draft very differently,
and a genuinely new product can be described in boilerplate
(Appendix~\ref{app:exhibit}). The measure sees prose.

\emph{Radicalness} \citep{DewarDutton1986, TushmanAnderson1986,
HendersonClark1990} is the closest neighbour, since KL is a departure
measure rather than a newness count and the Amazon exhibit is a large
departure built from unremarkable components. It is declined because it
claims that existing competences are devalued, and nothing here licenses the
step from unusual language to devalued competence.

\emph{Distinctiveness} has two prior occupants. In trademark law it is the
registrability doctrine, proven in part by the continued use this paper uses
as an outcome. In strategic management, optimal distinctiveness
\citep{Zhao2017} predicts an inverted-U in conformity --- a shape that
appears here at registration but not at listing, so the term would imply a
theory is being tested when it is not.

What \emph{atypicality} says is narrower than any of them: this description
is unusual for its class and year. Whether the thing described was new,
radical, or good is not measured.


